\documentclass[11pt]{article}
\usepackage[margin=2.5cm]{geometry}
\usepackage[T1]{fontenc}
\usepackage{lmodern}
\usepackage{amsmath,amssymb}
\usepackage{graphicx}
\usepackage{booktabs}
\usepackage{array}
\usepackage{xcolor}
\usepackage[numbers,sort&compress]{natbib}
\usepackage[colorlinks=true,linkcolor=blue!50!black,citecolor=blue!50!black,urlcolor=blue!50!black]{hyperref}
\usepackage{xurl}
\usepackage{authblk}
\usepackage{caption}
\usepackage[section]{placeins}
\captionsetup{font=small,labelfont=bf,labelsep=period}

% ----------------------------------------------------------------------------
\newcommand{\figfile}[2]{\IfFileExists{#1}{\includegraphics[width=#2]{#1}}{%
  \fbox{\parbox[c][4cm][c]{#2}{\centering\small missing file \texttt{\detokenize{#1}}}}}}
% ----------------------------------------------------------------------------

\title{A migration model of population recovery after disasters}

\author[1]{Kyuri Park}
\author[1]{Mike Lees}
\author[2]{Ella McCoshan}
\author[3]{Brendon McGuinness}
\affil[1]{Computational Science Lab, Informatics Institute, University of Amsterdam, Amsterdam, The Netherlands}
\affil[2]{Yale School of the Environment, Yale University, New Haven, CT, U.S.A.}
\affil[3]{Department of Biology and Department of Environmental Studies, New York University, New York, NY, U.S.A.}

\date{}

\setlength{\emergencystretch}{2em}
\AtBeginDocument{\lefthyphenmin=2\righthyphenmin=3}
% Set this switch to true to include the PNAS significance statement.
% The default follows the shorter Nature Cities article structure.
\newif\ifpnasversion\pnasversionfalse

\begin{document}



\maketitle

\begin{abstract}
Population recovery after disasters does not always follow the loss and rebuilding of housing. We test whether modelling departure, return, and settlement improves population predictions. A metropolitan migration model calibrated on Hurricane Katrina captures broad recovery patterns but misses some destinations and return timing. We then test a reduced county model across 21 counties affected by fourteen disasters. Each disaster is held out from fitting, and predictions use its observed housing recovery path. With departure and return triggered at the pre-disaster occupancy level, median population prediction error falls from 4.9 percentage points under direct capacity tracking to 3.6. A pooled occupancy threshold reduces it further to 3.0; this additional gain is smaller than the gain from modelling departure and return over time. The fixed-threshold model performs similarly to the strongest statistical forecasts, while the pooled model outperforms them. Even with only first-year housing loss and a typical rebuilding curve, the dynamic model is more accurate than capacity tracking. Performance varies across counties, and estimated thresholds change with the housing measure. For planning, the first post-event housing-loss estimate can be combined with a typical rebuilding curve to provide an initial population forecast, which can be updated as new housing estimates arrive. Direct measures of habitability and temporary accommodation would help explain remaining differences.
\end{abstract}

\ifpnasversion
\section*{Significance statement}
Housing reconstruction does not guarantee that residents will return after a disaster. We test a model that follows departure, return, and settlement across 21 counties affected by fourteen US disasters. Given the observed housing recovery, the model predicts population more accurately than assuming that population follows the housing stock. The fitted relation between population and capacity also changes with how housing is measured. This matters for recovery planning because temporary accommodation and uninhabitable homes can both alter the apparent relation between housing and population. The findings support modelling migration over time and collecting direct evidence on where affected residents stay.
\fi

\noindent\textbf{Keywords} urban recovery, housing loss, migration, metapopulation, natural disasters

After a disaster, the number of homes lost gives an incomplete picture of how many people will leave and when they will return. Some residents share homes or find temporary accommodation nearby. Others move to another city and remain there after housing is rebuilt. These movements affect demand for housing and services in both the affected area and the places receiving its residents \citep{Levine2007,Peacock2014}. Predicting recovery requires knowing how the population adjusts as housing is lost and rebuilt.

Research on Hurricane Katrina has documented several reasons why that adjustment differs across households. Return varied with age, income, race, and damage to the home \citep{Fussell2010,GroenPolivka2010}. County studies relate net migration after Katrina and Rita to dimensions of local social vulnerability \citep{Myers2008}. Many households moved within the affected region, while Houston and Baton Rouge received large numbers of residents \citep{SastryGregory2014,Hori2009}. Households also split and recombined after the storm \citep{Rendall2011}. Tax records show that some changes in where people lived persisted for years \citep{Deryugina2018}. Beyond Katrina, Kobe regained its pre-earthquake population within ten years, but residents shifted away from the older urban core \citep{Chang2010}. These findings make it unlikely that population recovery can be inferred directly from housing reconstruction.

Existing migration links help explain where people go. Movement after Katrina and Rita followed and strengthened connections between places \citep{Curtis2015,DeWaard2016}. Mobile phone records provide finer observations of departure and return \citep{Lu2012,Yabe2020}. Models of mobility and climate migration describe destination choices and possible changes in population distribution \citep{Simini2012,Hauer2017,Robinson2020}. These studies describe who returns and where displaced residents go. Whether population recovery can be predicted for a disaster not used to fit the model has received less attention. We ask whether a model that connects housing recovery to departure, return, and settlement can make such predictions.

We use a metapopulation framework, which links local populations through movement \citep{Levins1969,Hanski1998}. Housing capacity changes after a disaster. Population adjusts as residents leave, return, or settle elsewhere. This framework connects the conditions within an affected area to the places available to its residents, consistent with ecological approaches to urban systems \citep{Grimm2008,Seto2012,Espey2024}. The model does not assume that a rebuilt home immediately brings back its former occupants.

We examine the mismatch between housing and population recovery, develop a metropolitan network model for Katrina, and test a reduced county model across 21 counties affected by fourteen US disasters. For each prediction, every county affected by the target disaster is excluded from fitting, and the prediction uses the target county's observed housing path. The test therefore asks what migration dynamics add once housing loss and rebuilding are known. We compare the model with direct capacity tracking, a simpler dynamic specification, and statistical forecasts. We also examine whether the fitted occupancy threshold is stable across housing measures and whether hazard geography improves predictions. The pooled dynamic model predicts population more accurately than capacity tracking and the statistical forecasts. Most of its gain over capacity tracking comes from modelling departure and return over time.


\section*{Results}

\subsection*{Population recovery diverges from housing recovery}\label{sec:panel}

Population and housing recover at different rates after disasters. We first compare ten counties affected by Katrina, Rita, Harvey, and the Camp Fire. Seven of these counties were struck by Katrina. Across all ten, first post-event Census housing estimates record losses ranging from 3.5\% in Aransas County to 78\% in St Bernard Parish. We express both population and housing relative to their projected pre-disaster trends (Methods).

The differences run in both directions (Fig.~\ref{fig:events}). In Hancock, Harrison, and Butte, the population one year after the event exceeds the level implied by the remaining housing at pre-event occupancy by 6--20\%. This excess then narrows. In Hancock, it falls from 20\% in 2006 to 4\% in 2009. Temporary accommodation and shared housing could allow residents to remain locally. FEMA placed trailers and mobile homes across coastal Mississippi after Katrina \citep{GAO2007,EvansCowley2011}.

In Plaquemines, Cameron, and Aransas, the population falls below the level implied by the housing stock. In the first two counties, population continues to decline after housing stabilises. Damaged homes can remain in the housing count, and residents may leave for reasons that housing loss alone does not capture. These opposing patterns motivate a model that allows the population to adjust over time and treats housing capacity as an imperfectly measured constraint.

\begin{figure}[tbp]
\centering
\figfile{figures/fig1_events.pdf}{\textwidth}
\caption{Population and housing in ten affected counties, ordered by housing loss. Dots show Census population and squares show housing units, both relative to pre-event linear trends. The dotted curve is the capacity path fitted through the housing estimates. The dashed curve fixes the model's occupancy threshold at $\phi^*=1$. The solid curve fits a separate threshold for each county. Time is measured from landfall for Katrina, Rita, and Harvey, and from the Camp Fire for Butte.}
\label{fig:events}
\end{figure}

\subsection*{A model of departure, return, and settlement}\label{sec:calib}\label{sec:valid}

The model links local populations through migration (Fig.~\ref{fig:schematic}). A disaster reduces the affected area's housing capacity. Population is initially unchanged, so the number of residents relative to capacity rises. Residents then move towards areas with lower relative occupancy. Rebuilding allows residents to return. Some residents settle elsewhere and leave the group that may return, while receiving areas can add capacity.

We use two versions of this framework. The Katrina model represents a network of metropolitan areas. Migration links follow observed pre-disaster tax records, with a gravity model estimating suppressed flows. The county model used across disasters follows the affected county and its residents living elsewhere. It is the limit of the network model when receiving areas stay near their usual occupancy and local demographic change is omitted. With the Katrina rates, it reproduces the network's New Orleans population to within 0.4 percentage points (Methods). We use this reduced model across disasters. Tax records leave a median of 38\% of each county's out-migrants without a reported destination, so a reliable network cannot be built for each county.

In the county model, an occupancy threshold $\phi^*$ separates departure from return. When the population relative to capacity exceeds the threshold, residents leave; below it, residents who remain eligible to return can move back. A threshold of one corresponds to pre-event occupancy and recovers the network limit. When $\phi^*>1$, departure begins only when $N/K>\phi^*$; residents who remain eligible can return once $N/K<\phi^*$, even while population is still above measured capacity. When $\phi^*<1$, departure can begin while $N/K<1$, and return begins only once occupancy falls below $\phi^*$. The threshold can reflect local accommodation, other reasons to stay or leave, and errors in housing capacity.

The Katrina calibration illustrates why return needs its own dynamics. Tracking residents who may return fits New Orleans' population within 0.35 percentage points in each of the four annual observations. It also reproduces the decline in Census net domestic in-migration more closely than return driven only by the occupancy gradient. For the origins of excess in-migration in tax records, the variance-explained $R^2$ rises from 0.32 to 0.55 (Supplementary Fig.~\ref{fig:calib}).

The model captures broad destination patterns, with $R^2=0.70$, close to the 0.72 obtained from pre-disaster flow shares alone. Destination shares also informed the choice of coupling between receiving cities, so this comparison is not wholly independent of model selection. Important discrepancies remain. Houston receives 13\% of modelled departures and accounts for 13\% of returns, compared with observed shares of 28\% and 41\%. Return in the tax records is also more concentrated in the first year than the model predicts. Detailed calibration and validation results are in the Supplementary information.

\begin{figure}[tbp]
\centering
\figfile{figures/fig2_schematic.pdf}{\textwidth}
\caption{Model processes illustrated with rates calibrated on Hurricane Katrina. (a) The disaster reduces the affected city's housing capacity, shown in ochre. The dashed circle marks its original capacity, and blue shows population. Solid arrows indicate movement towards areas with lower relative occupancy. Dashed arrows indicate return as capacity recovers. (b) Population and capacity in the affected city over time. The shaded area is its contribution to cumulative excess occupancy, a simulation measure defined in the Supplementary information. (c) Population trajectories under different occupancy thresholds $\phi^*$. Values above one raise the occupancy level that triggers departure and allow return to begin while population is still above measured capacity. Values below one trigger departure even when population is below measured capacity, with return beginning only below the lower threshold. These thresholds describe aggregate responses rather than household preferences.}
\label{fig:schematic}
\end{figure}

\subsection*{Migration dynamics improve predictions across disasters}\label{sec:holdout}

We test population predictions in 21 counties affected by fourteen disasters. The sample combines the ten Census-panel counties with eleven additional counties whose American Community Survey (ACS) estimates show at least 4\% loss of usable housing. Usable housing excludes seasonal units and units vacant for other reasons. It remains a proxy for capacity rather than a direct measure of habitability. The sample and alternative inclusion thresholds are described in Methods.

For each disaster, we fit the model to all other disasters and predict the population in every county affected by the omitted event. One specification fixes $\phi^*=1$. The other estimates county thresholds around a common mean and uses that training mean for the omitted counties. Both use the omitted county's observed housing path. Its post-event population observations are excluded from fitting. These predictions are therefore conditional on measured housing loss and rebuilding.

The pooled threshold is 1.03 in the full fit and ranges from 1.02 to 1.03 across training sets. The pooled model has a median root-mean-square prediction error of 3.0 percentage points. Population proportional to housing capacity gives 4.9 points, and the dynamic model with $\phi^*=1$ gives 3.6 points (Fig.~\ref{fig:holdout} and Table~\ref{tab:heldout_summary}). The pooled model improves on capacity tracking in 16 of 21 counties. The median paired error reduction is 1.9 points, with a 95\% interval of 0.4--2.5 from a bootstrap that resamples whole disasters.

The additional benefit of pooling the threshold is smaller. Compared with the fixed-threshold dynamic model, the median paired improvement is 0.5 points, with a bootstrap interval of 0.2--1.3. A county-level Wilcoxon test gives $p=0.17$ and does not account for counties sharing a disaster. The fixed-threshold model itself improves on capacity tracking in 16 counties. Much of the gain, therefore, comes from allowing departure and return to unfold over time.

Statistical forecasts fitted to the same training disasters perform worse. Predicting the population gap from the capacity gap gives a median error of 5.6 points. Adding its interaction with time since the event gives 5.2 points, and adding first-year housing loss gives 4.0 points. A partial-adjustment forecast, in which the population gap moves towards the capacity gap at a fitted annual rate, gives 4.6 points. The dynamic model has lower error in 15 to 18 of the 21 counties in each comparison. Disaster-bootstrap intervals for the difference in median error favour the dynamic model in all four comparisons (Supplementary information). With $\phi^*=1$, the dynamic model also has lower error than each statistical forecast in most counties, but the intervals include zero for the three strongest forecasts.

Prediction remains uneven, with errors ranging from 0.8 to 10.4 points. The largest errors occur where the common threshold poorly represents the county, including Hancock, Cameron, Plaquemines, and Aransas. For the ten counties with at least 10\% housing loss, the median error is 3.6 points under the pooled model and 7.7 under capacity tracking. The median paired improvement is 2.2 points, but its interval spans $-0.2$ to 4.8. For small events, maintaining the pre-event population trend performs about as well as the dynamic model.

The main comparison persists under alternative pooling strengths and housing measures, but is less consistent when the sample changes. Raising the ACS inclusion threshold to 5\% leaves 17 counties and reduces the paired gain over capacity tracking to 1.4 points. Using ACS capacity for the five Census-panel counties with survey coverage leaves the pooled model's median error at 3.0 points. Yet errors increase in all four Katrina counties in this subset. Orleans rises from 4.6 to 13.3 points. Similar overall performance can conceal large changes in individual predictions.

Prediction intervals are broad. The 90\% intervals average 19 percentage points in width and cover 95\% of the 80 omitted county-year observations. They include variation in occupancy thresholds and evacuation rates while holding the housing path and shared parameters fixed.

Predictions remain useful when the rebuilding path is unknown. We repeated the prediction using only the first post-event housing estimate and the median rebuilding curve of the training disasters. The median error is then 2.5 points, compared with 4.5 when population follows the same capacity path. The dynamic model is more accurate in 19 of 21 counties. For the ten counties with at least 10\% housing loss, its median error rises from 3.6 to 4.2 points without the observed rebuilding path, compared with 8.3 for capacity tracking. Assuming no rebuilding gives 5.3 points for these counties. For smaller events, the typical curve slightly outperforms the measured path, probably because survey estimates of small housing losses are noisy (Supplementary information).

\begin{figure}[tbp]
\centering
\figfile{figures/fig3_holdout.pdf}{\textwidth}
\caption{Population prediction for disasters omitted from fitting. (a--h) The eight counties with the largest housing losses. Dots show Census population relative to trend and dotted curves show measured capacity. Predictions use parameters fitted to other disasters, with a pooled occupancy threshold shown by solid curves and $\phi^*=1$ by dashed curves. Both use the omitted county's observed capacity path. (i) Root-mean-square population error for all 21 counties, on a log scale and ordered by housing loss. Filled circles show the pooled model and squares show the $\phi^*=1$ model. Open circles assume that population follows capacity, while ticks assume no departure from trend. Medians appear above the panel in that order.}
\label{fig:holdout}
\end{figure}

\begin{table}[tbp]
\centering\small
\caption{Median county root-mean-square prediction error in percentage points. Each disaster is omitted from fitting in turn. All models use the same measured post-event capacity paths where applicable, and the statistical forecasts are fitted to the same training disasters. The second column includes all 21 counties. The third includes the ten counties that lost at least 10\% of housing.}
\label{tab:heldout_summary}
\begin{tabular}{@{}p{8.0cm}rr@{}}
\toprule
Model & All counties & Larger losses\\
\midrule
Dynamic model with pooled occupancy threshold & 3.0 & 3.6\\
Dynamic model with $\phi^*=1$ & 3.6 & 5.6\\
Regression on capacity gap, time, and first-year loss & 4.0 & 6.5\\
Partial adjustment towards measured capacity & 4.6 & 6.8\\
Population proportional to measured capacity & 4.9 & 7.7\\
Population remains on its pre-event trend & 6.3 & 13.2\\
\bottomrule
\end{tabular}
\end{table}

\subsection*{Housing measurement changes estimated occupancy thresholds}

Fitting a separate threshold to each of the original ten counties reduces the median population error from 2.8 to 0.5 percentage points (Fig.~\ref{fig:panel}). Estimated thresholds range from 0.80 to 1.25. This flexible model has 31 parameters for 47 observations. Its improved fit describes the population series used for estimation rather than demonstrating prediction for a new disaster.

The threshold depends strongly on the housing measure. Five of these counties also have ACS coverage. Replacing housing-unit counts with the usable-housing proxy raises the threshold in all five. Changes are modest in Harrison, Jackson, and Butte, but larger in Jefferson, from 1.01 to 1.13, and Orleans, from 0.97 to 1.27. With this measure, both New Orleans parishes also retain a population above the level implied by remaining usable housing. The Mississippi and Butte patterns persist across measures, while the interpretation of the New Orleans parishes changes.

Annual observations also limit what can be inferred about departure speed. Evacuation rates are identified in only five counties. Within that subset, the fitted rate falls as housing loss increases (Fig.~\ref{fig:panel}b). In the others, the population reaches the fitted threshold too quickly for annual data to distinguish evacuation rates. The apparent relation between shock size and departure speed rests on those five cases.

We find limited evidence for what explains threshold differences. Hazard geography is associated with fitted thresholds, especially within Katrina, but the association weakens when capacity is measured with the ACS proxy. Using the housing share outside mapped hazard zones to predict thresholds gives a median population error of 4.1 points across 17 storm and flood counties, compared with 3.7 for a common threshold. Elevation provides a modest improvement in a different coastal subset. However, the result is not consistent across measures (Supplementary Fig.~\ref{fig:hazard}).


Accommodation records cannot resolve what drives these threshold differences. FEMA hotel placements cover only six counties and do not show where other assisted households stayed. Household Pulse estimates count displaced people across whole states, whereas our capacity measures count housing units in counties. Because these sources measure different outcomes at different scales, $\phi^*$ is best read as an aggregate parameter that depends on how capacity is measured, not as evidence of a particular household response.

\begin{figure}[tbp]
\centering
\figfile{figures/fig4_panel.pdf}{\textwidth}
\caption{Fitted thresholds, evacuation rates, and population errors. (a) Occupancy threshold $\phi^*$ against first-year housing loss, with 95\% intervals. Values above one place the departure threshold above measured capacity and allow return to begin while population is still above capacity. Values below one place it below capacity, so departure can begin while measured capacity still exceeds population and return starts only below the threshold. The threshold also captures error in the capacity measure. St Bernard reaches the lower fitting bound, so its boundary-based interval is not meaningful. (b) Evacuation rate against housing loss for the five counties in which it is identified. Bars show 95\% intervals and the dashed line is the log--log fit. (c) Root-mean-square error of population relative to trend. Crosses use the New Orleans metropolitan rates. Open symbols fit each county's rates with $\phi^*=1$. Filled symbols also fit $\phi^*$.}
\label{fig:panel}\end{figure}

\section*{Discussion}

\subsection*{What migration adds to recovery estimates}



Housing loss and rebuilding help predict population recovery, but population does not simply follow the housing stock. Across the fourteen-disaster test, median prediction error is 4.9 percentage points when population is set proportional to measured capacity, 3.6 with the dynamic model's threshold fixed at one, and 3.0 with a pooled threshold. Most of the gain over capacity tracking comes from modelling departure and return over time; pooling the threshold adds a smaller median paired improvement of 0.5 points. The pooled model gives the clearest advantage over the strongest statistical forecasts. The fixed-threshold model has lower median error than the three strongest forecasts, but the disaster-level intervals do not distinguish it from them. The partial-adjustment forecast also changes gradually but has a higher median error than either dynamic model.

The Katrina comparison illustrates why return timing matters. In the model, spare housing alone does not determine return because some former residents have settled elsewhere. This interpretation is consistent with evidence that Katrina's effects on residence persisted for years \citep{Deryugina2018}. Our comparison does not identify the household decisions that produce settlement, but it shows the value of representing its consequences for population recovery.

For recovery planning, the first available housing-loss estimate can initialise a population forecast using a typical rebuilding curve from past disasters; the forecast can then be updated as new housing estimates arrive. This is a conditional forecast: before an event it requires a housing-loss scenario, and after an event it requires observed damage and rebuilding information. The dynamic model adds most over capacity tracking by representing departure and return, especially after larger housing losses. Rapid damage assessments would need to be calibrated against the first housing estimates in our data, which arrive within a year of each event. Our intervals hold the housing path and shared parameters fixed, so operational forecasts should also carry uncertainty in damage and rebuilding through to population estimates.

\subsection*{Interpreting occupancy thresholds}The threshold sets the occupancy ratio at which the model switches between outward movement and return. If $\phi^*>1$, outward movement occurs only when $N/K>\phi^*$, while eligible residents can return as soon as $N/K<\phi^*$, even if population still exceeds measured capacity. If $\phi^*<1$, outward movement can occur while measured capacity still exceeds population, and return starts only after $N/K$ falls below this lower threshold. Temporary housing and shared homes could contribute to a threshold above one, but the fitted value does not identify these mechanisms and also absorbs error in the housing-capacity measure.

Changing the housing measure changes some thresholds substantially. Census counts can include damaged homes that remain uninhabitable. ACS vacancy categories offer an alternative proxy, but they do not directly inspect habitability. The contrast between these measures is larger than the statistical uncertainty around several fitted thresholds. The threshold should therefore not be read as a measured ability to accommodate residents.

Neither hazard geography nor the small samples of assistance and survey records resolve this ambiguity. Direct records of temporary accommodation, housing habitability, and moves within the county would make these explanations testable.

\subsection*{Limits of the evidence and next steps}

The sample is constrained by the housing data. Census estimates miss losses after some disasters. Annual ACS estimates exclude smaller counties and have sampling errors that can be large relative to the measured losses. Our fourteen-disaster comparison combines two capacity measures and includes eleven additional ACS counties selected for a recorded loss of at least 4\%. Sensitivity analyses help assess these choices, but the sample does not represent all disasters or all affected communities.

The prediction test excludes whole disasters, but only fourteen events are available. Several evacuation rates are weakly identified from annual observations. Large prediction errors also remain where a common occupancy threshold poorly describes a county. Better measurements of capacity and household moves would help distinguish these errors from inadequate migration rules.

The network model has further limits. It underestimates Houston's role in Katrina and misses the timing of return in the tax records. Organised evacuation is a plausible reason for the destination mismatch \citep{McIntosh2008}. IRS data also omit non-filing households and suppress small flows. Return rules do not distinguish households by tenure, race, or income, despite their relevance after Katrina \citep{Fussell2010,GroenPolivka2010}. Future extensions should be tested on disasters excluded from fitting, using the same housing information and comparison models.

Receiving areas' housing response also needs direct measurement. In the calibrated network simulations, faster capacity expansion reduces cumulative excess occupancy, but the response rate is poorly constrained by the current housing targets. Building permits, vacancy data, and temporary housing records could test this prediction. Together with household histories, these observations would show when a migration model improves population forecasts and when a more detailed account of accommodation and return is needed.


\section*{Methods}\label{sec:methods}\label{sec:model}

\subsection*{Data and spatial units}

The Katrina model links 196 metropolitan areas, defined using the 2009 Census delineation. Each has at least 250{,}000 residents in 2005 or lies within 600~km of New Orleans. New Orleans denotes the seven-parish metropolitan area. Comparisons across disasters use counties. These two spatial scales serve different analyses and are not treated as interchangeable.

IRS Statistics of Income migration files record changes in filing addresses between consecutive years. Exemptions approximate the people represented by tax returns. We use 2003--04 to 2007--08 files for Katrina, 2016--17 to 2020--21 for the Camp Fire, and county totals extending to 2022--23 for the larger panel. For county-to-county flows, cells with fewer than ten returns were suppressed through 2012--13 and fewer than twenty from 2013--14 onward. US-total county records include these suppressed flows. The data exclude non-filers and do not observe all household moves. Filing years also differ from calendar years.

For example, the 2005--06 file compares addresses on returns filed in consecutive springs and covers approximately April 2005 to April 2006. We aggregate county flows to metropolitan areas before constructing network weights. The metropolitan sample contains 222 million residents and receives 88\% of destination-identified flows out of New Orleans in 2005--06. These coverage figures do not include destinations hidden by disclosure suppression.

Census intercensal estimates for 2000--2009 and Vintage 2020 estimates provide population and housing stocks. Vintage 2009 metropolitan components of change provide natural increase and net domestic migration. ACS 1-year tables B01003, B25001, B25002, B25004, B25010, and B11001 provide population, housing, vacancy, and household measures for 2005--2024. There is no standard 2020 ACS release. Annual ACS coverage excludes counties with fewer than 65{,}000 residents. Source files, checksums, and processing scripts are provided with the code.

\subsection*{Network population dynamics}

Each metropolitan area has population $N_i(t)$ and housing capacity $K_i(t)$, measured in people. Capacity is the number of residents the housing stock can accommodate at its reference occupancy. Relative occupancy is $\phi_i=N_i/K_i$. A value above one indicates population above that reference. It does not establish that every household is crowded.

Before the disaster, $N_i=K_i$. Local growth and occupancy-driven migration are then zero. The model describes changes around this baseline. Time is measured in years. Population changes through local demographic adjustment, movement along the network, and return migration,
\begin{equation}
\begin{aligned}
\frac{dN_i}{dt} &= r_i N_i (1-\phi_i) + \sum_j\left(F_{ji}-F_{ij}\right) + R_i, \\
F_{ij} &= W_{ij}\,N_i\,(\phi_i-\phi_j)_+ .
\end{aligned}
\label{eq:N}
\end{equation}
Here $(x)_+=\max(x,0)$. Movement occurs toward lower relative occupancy. The connection weight $W_{ij}$ has units of inverse years. The logistic term is informed by natural increase and becomes negative above capacity. It represents net local population adjustment rather than observed deaths. The term $R_i$ records return migration. Every migration flow is subtracted from its origin and added to its destination, so migration conserves total population.

\subsection*{Housing loss and rebuilding}

The disaster occurs at $t=0$ in area $s$. It removes a fraction $d_0$ of the original capacity. A fraction $d_\infty$ remains lost in the long run, while the recoverable part is rebuilt at rate $\rho_s$,
\begin{equation}
K_s(t)=K_s(0^-)\left[1-d_\infty-(d_0-d_\infty)e^{-\rho_s t}\right],\qquad t\geq0.
\label{eq:Ks}
\end{equation}
Population is unchanged at the instant of the shock. The resulting increase in occupancy creates pressure to leave. We scale outgoing weights so that $\sum_jW_{sj}=S_s$. If all destinations remain at occupancy one, total departure is $S_sN_s(\phi_s-1)_+$. In the full network it also depends on destination occupancy.

Receiving areas add capacity when population exceeds it,
\begin{equation}
\frac{dK_i}{dt}=\rho_h(N_i-K_i)_+,\qquad i\ne s.
\label{eq:hosts}
\end{equation}
The parameter $\rho_h$ controls the response speed. This rule produces no capacity growth below the reference occupancy. Host-to-host movement is scaled by $S_{\text{med}}$, the median sum of outgoing host weights. Its fitted value is zero. This choice concerns aggregate net flows and does not exclude onward moves by individual residents.

\subsection*{Return and settlement}

The variable $D_j$ tracks residents from the affected area who live in host $j$ and may return. They are already included in $N_j$. They return when the affected area has spare capacity and leave the return-eligible group as they settle elsewhere,
\begin{equation}
\begin{aligned}
\frac{dD_j}{dt} &= F_{sj}-\kappa(1-\phi_s)_+D_j-\lambda D_j, \\
R_j &= -\kappa(1-\phi_s)_+D_j, \\
R_s &= \sum_j\kappa(1-\phi_s)_+D_j .
\end{aligned}
\label{eq:D}
\end{equation}
The parameters $\kappa$ and $\lambda$ control return and settlement. Settlement reduces $D_j$ but leaves those residents in the host population. Return reduces the host population and increases the affected population equally. Ordinary network flows into $s$ are disabled to avoid counting return twice. A comparison model omits $D_j$ and instead lets return follow the occupancy gradient along incoming connections, scaled by $S_{\text{ret}}$.

\subsection*{Migration connections}

Network weights are proportional to observed pre-disaster IRS flows. We estimate suppressed connections with a gravity model,
\begin{equation}
w_{ij}\propto P_i^\alpha P_j^\beta d_{ij}^{-\gamma},
\label{eq:gravity}
\end{equation}
where $P_i$ and $P_j$ are populations and $d_{ij}$ is great-circle distance between Census internal points. We fit mean 2003--04 and 2004--05 flows across all 38{,}220 ordered metropolitan pairs using Poisson pseudo-maximum likelihood \citep{SantosSilva2006}. Estimates use HC1 standard errors and a metropolitan cluster bootstrap with 200 replicates. Imputed flows are capped at 20 exemptions, approximating the disclosure threshold. Reported flows retain their observed weights. Gravity supplies missing connections rather than the dynamic rule for migration.

Excess migration is the observed flow minus its pre-disaster mean. Destination-share analyses set negative excess flows to zero. Supplementary Methods describe the New Orleans-specific gravity estimates and the separate comparison of destinations after seven other county events.

\subsection*{Katrina calibration and validation}

July 2005 populations set initial capacities. Rates of natural increase in 2002--05 set $r_i$. Stage A fits Eq.~\eqref{eq:Ks} to metropolitan housing estimates for 2006--09, divided by the projected 2001--05 linear trend. Stage B fits $(S_s,\kappa,\lambda)$ to population relative to trend in those years and the share of two-year excess out-migration recorded in 2005--06. Stage C fits $\rho_h$ to housing growth above trend in 2006 and 2007 across the ten main receiving metropolitan areas. Housing increments are converted to capacity using each area's 2005 persons per unit. Four metropolitan areas directly affected by Katrina or Rita are excluded from receiving-area targets and validation.

We solve the equations with LSODA at relative tolerance $10^{-6}$. Stage A uses nonlinear least squares and a leave-one-year-out jackknife. Stages B and C use weighted nonlinear least squares on log parameters. Residual variances are estimated and updated once. We alternate these stages four times. Parameter covariance uses the Gauss--Newton approximation.

The initial capacity loss is extrapolated from annual housing observations rather than measured at landfall. The first observation in July 2006 therefore affects how the fit separates initial damage from the evacuation rate. To assess this dependence, we repeat stages B and C after omitting each year from the stage A housing fit.

For the Katrina uncertainty bands, we take 200 draws of $(S_s,\kappa,\lambda)$ from their joint log-scale covariance and draw $\rho_h$ from the lognormal distribution represented by its interval. These bands describe uncertainty in the calibrated rates. They do not include all uncertainty in the housing observations or initial capacity loss.

Flows are accumulated to April for IRS comparisons and to 30 June for Census components of change. Additional comparisons use destination shares, return origins and timing, receiving-area net migration, and occupancy. These were not targets for the evacuation, return, and settlement rates. Destination shares did inform the host-to-host coupling choice. Uncertainty propagation and alternative network definitions are detailed in Supplementary Methods.

\subsection*{County samples and capacity measures}

The first panel contains ten counties with housing losses recorded in Census unit estimates. Seven were affected by Katrina, with one each affected by Rita, Harvey, and the Camp Fire. Population and housing are divided by linear trends fitted to five pre-event July estimates. Follow-up ends in July 2009 for the 2005 disasters and July 2020 for the later events. With at least four post-event housing observations, capacity follows Eq.~\eqref{eq:Ks} fitted to housing ratios. Otherwise it falls initially to the first observed ratio and changes linearly between later observations.

The Katrina counties are Orleans, St Bernard, Plaquemines, and Jefferson in Louisiana, and Hancock, Harrison, and Jackson in Mississippi. The other cases are Cameron after Rita, Aransas after Harvey, and Butte after the Camp Fire. Their pre-event trend windows are 2001--05, 2013--17, and 2014--18 for the 2005 disasters, Harvey, and the Camp Fire, respectively.

The expanded analysis uses ACS usable housing, defined as total units minus seasonal units and units vacant for other reasons. The latter category includes some damaged homes. This measure remains a proxy for habitability. Its reference trend uses at least three of the six ACS years before the event, excluding the event year and placing observations at mid-year.

We listed 42 county--disaster pairs before fitting and added an omitted case, Orleans after Ida, giving 43. Thirty-seven have sufficient ACS observations. Eleven additional counties have first-year usable-housing losses of at least 4\%. Together with the original ten, they form the prediction sample of 21 distinct counties across fourteen disasters. Butte after the Camp Fire is counted once. Sensitivity analyses use 3\% and 5\% inclusion thresholds.

Expanded-panel population trends use five pre-event July estimates. We combine intercensal estimates, Vintage 2020, and Vintage 2025, rescaling the latter to match Vintage 2020 at 2020. Follow-up lasts 4.6 years or ends at another disaster. The ACS capacity check for Katrina uses a Census-based reference because pre-event annual ACS data are unavailable. Its construction is given in Supplementary Methods.

For Orleans, Jefferson, Harrison, and Jackson, that reference combines the 2000--05 Census housing trend with the Census 2000 share of usable units. ACS usable housing is divided by the resulting series. Butte uses the expanded-panel ACS measure. We refit each of these counties separately, holding $\kappa$ at its ten-county estimate and allowing $S_s$, $\lambda$ and $\phi^*$ to vary. Multiple starting points reduce dependence on the initial parameter values.

\subsection*{Reduced county model}



To reduce the network to one affected county, hold every receiving area's occupancy at its reference value ($\phi_j=1$) and omit local demographic change. Summing the origin flows in Eq.~\eqref{eq:N} then gives total departure $E=S_sN_s(\phi_s-1)_+$. Summing Eq.~\eqref{eq:D} across hosts, with $D=\sum_jD_j$, gives return $A=\kappa(1-\phi_s)_+D$ and settlement $\lambda D$. The affected population therefore changes as $\dot N_s=-E+A$, while the return-eligible pool changes as $\dot D=E-A-\lambda D$. These are the county equations in Eq.~\eqref{eq:county} when $\phi^*=1$, before expressing population and capacity relative to their projected trends. The fitted county model generalises this limiting case by replacing one with $\phi^*$ in both departure and return. The reduction is close for Katrina: no receiving area's occupancy in the calibrated network rises more than 3.9\% above its reference, and the excess occupancy averaged over displaced residents stays below 1\%. Using the same rates and $\phi^*=1$, the county equation reproduces the network's New Orleans population to within 0.36 percentage points over four years.

We use the county model for the cross-disaster test because tax records do not support a destination network for each county. In the filing year before each event, individually reported county-to-county flows account for a median of 62\% of out-migrants (range 29--87\%, for the 19 of 21 counties with available files). The smallest counties report as few as two destinations. The remaining flows are suppressed into aggregates without a destination. The implementation of both checks is in \texttt{28\_reduction\_check.py}.

In the county model, $N(t)$ and $K(t)$ are population and housing ratios relative to their separate projected trends. We fit the model directly to these ratios, starting from $N(0)=1$ and $D(0)=0$. The unobserved pool $D$ represents residents who may return on the same normalised model scale.

An occupancy threshold $\phi^*$ separates departure from return,
\begin{equation}
\begin{aligned}
E(t)&=S_sN\left(N/K-\phi^*\right)_+, \\
A(t)&=\kappa\left(\phi^*-N/K\right)_+D, \\
\frac{dN}{dt}&=-E(t)+A(t), \\
\frac{dD}{dt}&=E(t)-A(t)-\lambda D.
\end{aligned}
\label{eq:county}\end{equation}

The threshold shifts the movement trigger relative to measured occupancy. If $\phi^*>1$, departure occurs when $N/K>\phi^*$, and eligible residents can return when $N/K<\phi^*$, even while population remains above capacity. If $\phi^*<1$, departure can occur while $N/K<1$, and return begins only once occupancy falls below $\phi^*$. The fitted threshold can reflect local accommodation, other reasons to stay or leave, and capacity measurement error. It is not a direct measure of household preferences or a separately identified accommodation mechanism.

Fits minimise weighted squared errors in population ratios and the event-year share of two-year excess IRS out-migration. Measurement standard deviations are fixed at 0.005 and 0.03 respectively. The timing share is $e_1/(e_1+\max(e_2,0))$, where excess flows use the two preceding years as baseline. In the ten-county panel, $\kappa$ is shared and $(S_s,\lambda)$ are county-specific. Extended fits also estimate county thresholds within $[0.8,1.4]$. The expanded panel shares both $\kappa$ and $\lambda$, with thresholds bounded by $[0.7,1.5]$. Supplementary Methods describe profile likelihoods and alternative capacity fits.

Gauss--Newton intervals use the assumed measurement errors. Profile likelihoods fix $\phi^*$ on a grid and re-optimise evacuation and settlement rates. A narrow interval under one capacity measure does not include uncertainty about that measure. Annual observations also constrain evacuation rates poorly when population reaches the fitted threshold between observations.

\subsection*{Prediction for omitted disasters}

We compare two dynamic models using the 21-county sample. One fixes $\phi^*=1$. The other partially pools county thresholds by adding $\sum_j(\phi_j^*-\mu)^2/\tau^2$ to the weighted fitting loss. We estimate $\mu$ and set $\tau=0.1$. Both models share return and settlement rates across counties. Evacuation rates vary by county, as do thresholds in the pooled model.

For each disaster, we omit all affected counties and refit to the remaining disasters. We regress fitted $\log S_s$ on log housing loss within the training sample. Each omitted county receives the evacuation rate predicted by this relation. Its threshold is the training mean $\mu$ in the pooled model or one in the fixed-threshold model.

Predictions use the omitted county's measured post-event housing path. Its post-event population data are excluded from fitting. The test therefore predicts population conditional on observed housing loss and rebuilding. It does not forecast housing damage or recovery from information available before the event. We report each county's root-mean-square population error in percentage points and compare errors across the same omitted-disaster groups.

A second test uses less housing information. The omitted county's capacity becomes $1-L_1m(t)/m(t_1)$. Here $L_1$ is its first post-event capacity loss, observed at time $t_1$, and $m(t)$ is the median across training counties of capacity loss divided by each county's own first loss. A variant holds capacity at $1-L_1$. The first post-event estimates come 2--11 months after the event, with a median of 10 months. The implementation is in \texttt{27\_early\_damage.py}.

\subsection*{Benchmarks and uncertainty}

Two simple benchmarks set population equal to measured capacity or maintain its pre-event trend. Three additional regressions use the same training disasters and have no intercept. The first fits $N-1=b(K-1)$. The second adds $(K-1)t$, and the third also adds the first-year capacity gap $K_0-1$. A partial-adjustment forecast starts at the pre-event trend and moves the population gap $x$ towards the capacity gap $g$ at each observation time, $x_k=x_{k-1}+(1-e^{-a\Delta t_k})(g_k-x_{k-1})$, with the rate $a$ fitted to the training disasters. It is computed recursively and uses no post-event population observation of the omitted county.

Uncertainty in model comparisons uses 10{,}000 bootstrap draws of whole disasters (2{,}000 for the sensitivity variants), keeping affected counties together. We also report county-level paired Wilcoxon tests, which do not account for that clustering. Sensitivity checks vary $\tau$ to 0.05 or 0.2, the ACS inclusion threshold, and the housing measure in counties with both sources. These checks and the Katrina sensitivity analysis are detailed in Supplementary Methods and Supplementary Section~\ref{si:sensitivity}.

The difference between models' median errors and the median of their paired county differences are distinct summaries. We report paired error reductions for the main capacity-tracking comparison. Bootstrap intervals for the statistical forecast comparisons describe differences in median error. Counties remain weighted equally when calculating these summaries, while disaster-level resampling accounts for shared events.

Prediction intervals use 400 draws per county. Thresholds are drawn from $N(\mu,0.1^2)$ and clipped to $[0.7,1.5]$. Log evacuation rates are drawn around the training regression using its residual standard deviation. The observed housing path, shared rates, and regression coefficients remain fixed. The intervals therefore represent selected between-county variation and exclude uncertainty in those fixed inputs.

\subsection*{Geography and accommodation checks}

For seventeen storm- and flood-affected counties, we calculate the share of pre-event housing outside mapped hazard areas and, for coastal counties, above 3~m elevation. Hazard areas combine current FEMA flood zones, NOAA Category~3 surge zones, and NOAA levee areas. Inland flood cases use flood zones alone. Definitions were fixed before calculating full-sample shares, after a preliminary Katrina check. Current maps make this a retrospective comparison. Geography enters the pooled threshold mean as $a+bx_j$, estimated within each training sample with $\tau=0.1$.

We locate the internal point of each 2010 Census block and weight hazard and elevation classifications by housing units. Events from 2010 onward use 2010 block housing. For earlier events, Census 2000 block housing is allocated to 2010 tracts by intersecting land area, then distributed within each tract in proportion to 2010 housing. This assumes the later spatial distribution within tracts also represents the earlier period. Comparable fire-hazard maps are not included.

Separate checks use OpenFEMA registrations and hotel placements, and Household Pulse reports of disaster displacement. These sources differ in geography, coverage, and timing. They cannot directly identify where all affected households stayed. Supplementary Methods give spatial weighting, survey variance estimation, and source-specific definitions. The simulation measure of cumulative excess occupancy is defined separately in Supplementary Section~\ref{si:sim_measure}.

FEMA hotel stays count registrants who checked in under Transitional Sheltering Assistance. Rental-resource location fields are almost always blank. Household Pulse estimates use person weights and 80 replicate weights, with successive-difference variances scaled by $4/80$. The comparison with adults associated with lost housing multiplies ACS usable units lost by persons per household and an adult share of 0.78. State survey estimates and county housing totals have different geographic coverage.

\section*{Data and code availability}\label{sec:availability}
IRS SOI migration data are available at \url{https://www.irs.gov/statistics/soi-tax-stats-migration-data}. Census population and housing estimates and Gazetteer files are available at \url{https://www2.census.gov}. ACS tables are accessed through the Census Data API at \url{https://api.census.gov}, using a free API key. FEMA Individual Assistance registrations and registrant records are available through OpenFEMA at \url{https://www.fema.gov/about/openfema/data-sets}. Household Pulse Survey public-use files are at \url{https://www.census.gov/programs-surveys/household-pulse-survey.html}. Census blocks were obtained from TIGERweb (\url{https://tigerweb.geo.census.gov}); flood hazard areas from the Esri Living Atlas copy of the FEMA National Flood Hazard Layer; storm-surge hazard maps from the NOAA National Hurricane Center; and elevation from Tilezen terrain tiles. The hazard classification computed from these services is included in the repository.

The code, processed data, raw-file checksums, tests, and model audit are in \url{https://github.com/KyuriP/cities-project}, in the \texttt{empirical/} folder. The script \texttt{run\_all.sh} reproduces the reported numbers and figures from the raw files.

\section*{Acknowledgements}
This work began as a group project at the Santa Fe Institute Complex Systems Summer School 2026.

\bibliographystyle{unsrtnat}
\bibliography{refs}

\clearpage
\appendix
\renewcommand{\thetable}{S\arabic{table}}
\renewcommand{\thefigure}{S\arabic{figure}}
\renewcommand{\thesubsection}{S\arabic{subsection}}
\setcounter{table}{0}
\setcounter{figure}{0}
\setcounter{subsection}{0}
\section*{Supplementary information}

\subsection{Hurricane Katrina in the data}\label{si:data}\label{sec:data}
\subsubsection{Sources}
We combine public migration, population, and housing data from the IRS and the Census Bureau (Methods). The IRS county-to-county migration files count tax returns and exemptions whose filing address changed between consecutive filing years. Exemptions approximate the number of people represented by those returns. We use filing years 2003--04 to 2007--08, with about 80{,}000 reported county pairs per year.

Census estimates give each county's population and housing stock on 1 July. For metropolitan areas, the Census also reports births, deaths, and net domestic and international migration for years ending 30 June. We aggregate counties into the 366 metropolitan areas in the 2009 delineation. New Orleans refers to the seven-parish metropolitan area.

The model includes 196 metropolitan areas. Each has at least 250{,}000 residents in 2005 or lies within 600~km of New Orleans. Together they contain 222 million people and receive 88\% of the destination-identified flows out of New Orleans in 2005--06. Download and processing scripts record file checksums and compare their outputs with independently extracted values.

The IRS data omit households that do not file tax returns. These households are disproportionately poor and were disproportionately represented among Katrina evacuees \citep{Fussell2010}. Flows involving fewer than ten returns are suppressed. Only 31--32\% of normal out-migration from New Orleans can therefore be assigned to a destination county, compared with 71\% in 2005--06. Filing years also differ from calendar years. The 2005--06 file compares addresses on returns filed in consecutive springs and covers roughly April 2005 to April 2006.

Census population and housing figures are estimates. After Katrina, the Bureau produced special housing estimates for the affected parishes. Other disasters also affected potential receiving areas. Rita struck Beaumont and Lake Charles four weeks after Katrina, while Katrina itself struck the Mississippi coast. We exclude these four metropolitan areas from targets and validation involving receiving cities.

\subsubsection{Population and housing after Katrina}
\begin{figure}[tbp]
\centering
\figfile{figures/figS1_map.pdf}{\textwidth}
\caption{Migration from and to New Orleans after Hurricane Katrina. (a) Excess out-migration in filing year 2005--06, by destination metropolitan area. Excess is the flow above its 2003--05 mean. (b) Excess in-migration in 2006--07, by origin metropolitan area. Circle area and line width represent IRS exemptions, aggregated from county flows. Suppressed flows are omitted. Grey dots mark the 196 metropolitan areas in the network. The star marks New Orleans.}
\label{fig:map}
\end{figure}

Orleans Parish lost 53\% of its residents between July 2005 and July 2006. Across the wider metropolitan area, population fell to 0.740 of its 2001--05 trend in July 2006 and recovered to 0.813 by 2009. Housing recovered faster, rising from 0.727 to 0.871 of trend over the same period. Net domestic migration was $-332{,}206$ in the year to June 2006. Its 2002--05 average had been $-6{,}944$.

The IRS data show the same abrupt outward movement. Annual out-migration from the metropolitan area rose from about 62{,}000--64{,}000 before the storm to 258{,}823 in 2005--06. The increase above baseline was 195{,}793. Excess in-migration peaked one filing year later, at 32{,}755 in 2006--07.

The destinations also changed (Fig.~\ref{fig:map}). Only 26\% of excess out-migration remained within 300~km, compared with 45\% of baseline out-migration. Houston received 28\% of the excess movement bound for metropolitan areas. It was also the origin of 41\% of excess return migration. Houston's net domestic migration in the year to June 2006 was 82{,}043 above its 2002--05 average. By July 2006, its housing stock was 14{,}071 units above trend.

Table~\ref{tab:desc} and Fig.~\ref{fig:data} provide the underlying migration, population, and housing comparisons.

\begin{table}[tbp]
\centering\small
\caption{Descriptive statistics for the New Orleans metropolitan area. The baseline is the mean of 2003--04 and 2004--05. Identified flows are those assigned to a county. Concentration measures use the 195 receiving metropolitan areas.}
\label{tab:desc}
\begin{tabular}{lrrrrr}
\toprule
\multicolumn{6}{l}{\emph{A. IRS migration, by filing year}}\\
& 2003--04 & 2004--05 & 2005--06 & 2006--07 & 2007--08\\
\midrule
Out-migration, all & 61{,}718 & 64{,}343 & 258{,}823 & 71{,}891 & 62{,}511\\
In-migration, all & 58{,}507 & 58{,}647 & 69{,}189 & 91{,}332 & 76{,}087\\
Share identified by county & 0.31 & 0.32 & 0.71 & 0.43 & 0.34\\
\midrule
\multicolumn{6}{l}{\emph{B. Geography of flows}}\\
& \multicolumn{2}{r}{Baseline out} & Excess out & Return in & \\
& \multicolumn{2}{r}{2003--05} & 2005--06 & 2006--07 & \\
\midrule
Share within 300 km & \multicolumn{2}{r}{0.45} & 0.26 & 0.22 & \\
Largest metro share & \multicolumn{2}{r}{0.14 (Baton Rouge)} & 0.28 (Houston) & \multicolumn{2}{l}{0.41 (Houston)}\\
Herfindahl index & \multicolumn{2}{r}{0.060} & 0.122 & 0.208 & \\
\midrule
\multicolumn{6}{l}{\emph{C. Census, New Orleans metro relative to 2001--05 trend (1 July)}}\\
& & 2006 & 2007 & 2008 & 2009\\
\midrule
Population & & 0.740 & 0.774 & 0.797 & 0.813\\
Housing units & & 0.727 & 0.812 & 0.857 & 0.871\\
\bottomrule
\end{tabular}
\end{table}

\begin{figure}[tbp]
\centering
\figfile{figures/figS2_data.pdf}{0.95\textwidth}
\caption{Hurricane Katrina in the data. (a) Net domestic migration in New Orleans (black) and the four largest receiving metropolitan areas (grey). Census components of change cover years ending 30 June. (b) Population (dots) and housing units (squares) in the New Orleans metropolitan area on 1 July, relative to their 2001--05 linear trends. (c) Cumulative share of migrants by distance from New Orleans, using IRS county flows. (d) Shares of migrants moving to metropolitan areas received by the ten largest destinations of the 2005--06 excess. Open symbols show the baseline shares, and filled symbols show the storm-year shares.}
\label{fig:data}
\end{figure}

\subsection{Network estimation and Katrina calibration}\label{si:calibration}
\subsubsection{Fitting the model in stages}
We fit separate groups of parameters to housing, population, and migration data (Table~\ref{tab:calib}). July 2005 populations set the pre-disaster capacities. The 2002--05 rates of natural increase set $r_i$. Supplementary Section~\ref{si:validation} compares the fitted model with additional outcomes.

Stage A fits the New Orleans capacity path in Eq.~\eqref{eq:Ks} to metropolitan housing estimates for 2006--09, expressed relative to trend. Stage B fits the evacuation, return, and settling rates $(S_s,\kappa,\lambda)$. Its targets are population relative to trend in 2006--09 and the timing of outward movement. For the latter, we use the share of two-year excess IRS out-migration recorded in 2005--06, which is 0.957.

Stage C estimates the capacity response $\rho_h$ in receiving cities. We use housing growth above trend in 2006 and 2007 for the ten metropolitan areas that received the most evacuees. Each housing increment is converted to population capacity using that area's 2005 persons per housing unit. We alternate stages B and C until the estimates stabilise (Supplementary Section~\ref{si:methods}). Residual variances are estimated from each fit. Standard errors use the Gauss--Newton approximation on the log scale.

\begin{table}[tbp]
\centering\small
\caption{Calibrated parameters and their data sources. Intervals are 95\% Gauss--Newton intervals unless a standard error or another method is specified.}
\label{tab:calib}
\begin{tabular}{@{}l>{\raggedright\arraybackslash}p{4.1cm}>{\raggedright\arraybackslash}p{3.6cm}>{\raggedright\arraybackslash}p{4.4cm}@{}}
\toprule
Parameter & Meaning & Estimate & Determined by\\
\midrule
$K_i(0^-)$, $r_i$ & initial capacity and growth & 2005 population and natural increase & Census\\
$\alpha,\beta,\gamma$ & gravity estimates for suppressed pairs & 1.02, 0.90, 1.12 & IRS, 38{,}220 metro pairs, 2003--05\\
$d_0$ & capacity lost at landfall & 0.42 (jackknife SE 0.16) & Census housing units 2006--09\\
$d_\infty$ & capacity lost permanently & 0.11 (SE 0.02) & Census housing units 2006--09\\
$\rho_s$ & rebuilding rate & 0.76 yr$^{-1}$ (SE 0.30) & Census housing units 2006--09\\
$S_s$ & evacuation rate & 1.62 yr$^{-1}$ (1.54--1.69) & Census population 2006--09 and IRS timing\\
$\kappa$ & return rate per unit spare capacity & 8.8 yr$^{-1}$ (6.1--12.5) & Census population 2006--09 and IRS timing\\
$\lambda$ & settling rate & 0.50 yr$^{-1}$ (0.40--0.64) & Census population 2006--09 and IRS timing\\
$\rho_h$ & receiving-city capacity response & 1.9 yr$^{-1}$ (0.4--8.6) & receiving-city housing units 2006--07\\
$S_{\text{med}}$ & coupling between receiving cities & 0 & IRS onward flows, with no excess detected\\
\bottomrule
\end{tabular}
\end{table}

\subsubsection{Estimating the migration network}
We estimate the gravity model from pre-disaster flows using Poisson pseudo-maximum likelihood \citep{SantosSilva2006}. The fit includes all 38{,}220 ordered pairs of the 196 metropolitan areas. Only 38\% of pairs have a reported flow. The estimates are $\alpha=1.02\pm0.04$, $\beta=0.90\pm0.03$ and $\gamma=1.12\pm0.03$. A bootstrap that resamples metropolitan areas gives intervals of 0.89--1.20, 0.77--1.07 and 0.94--1.34, respectively.

New Orleans has a steeper distance gradient than the national network. At county level, its estimated distance exponent is $\gamma=2.16\pm0.09$ (Supplementary Table~\ref{tab:gravity}). This difference motivates using observed flows wherever they are available. Gravity estimates fill only the suppressed pairs. As a static predictor of evacuee destinations across receiving metropolitan areas, the national gravity model gives $R^2=0.35$. A New Orleans-specific gravity model gives 0.62, while observed pre-disaster flow shares give 0.72.

The estimated distance effect changed little after Katrina. The county-level exponent was $2.16\pm0.09$ for baseline out-migration and $2.07\pm0.14$ for excess out-migration after the storm. The destination-size exponent increased from $1.42\pm0.06$ to $1.57\pm0.15$ at county level. At metropolitan level, it increased from $1.41\pm0.07$ to $1.66\pm0.20$. These estimates are consistent with longer moves toward large cities despite a similar distance effect. Return migration in 2006--07 followed a similar pattern.

We also evaluate seven affected counties with released IRS flows, covering the Camp Fire, Michael, Laura, Irma, Florence, and two counties struck by Harvey. In this analysis, performance is measured by the squared correlation between predicted and observed destination shares. Its median is 0.17 for pre-disaster flow shares, compared with 0.03 and 0.06 for the gravity alternatives. Performance is strongest for the three largest outflows, with squared correlations of 0.44 in Butte, 0.28 in Calcasieu and 0.17 in Bay. These scores use a different definition from the variance-explained $R^2$ reported for Katrina. In four of the seven counties, excess out-migrants travelled farther than baseline migrants (Supplementary Section~\ref{si:methods}).

\subsubsection{Movement between receiving cities}
\label{si:host}
We find no detectable excess net out-migration from receiving cities to other metropolitan areas, including cities with large evacuee inflows. Fitting the coupling between receiving cities to destination shares also places its estimate at zero. We therefore set $S_{\text{med}}=0$. This result concerns aggregate net flows. It does not rule out onward moves by individual evacuees.

We look for onward migration by comparing receiving cities' excess net flows to one another with the number of evacuees they received. The regression slopes are $-0.39\pm0.16$ in 2005--06, $-0.64\pm0.20$ in 2006--07, and $-0.61\pm0.10$ in 2007--08. The large Sunbelt metropolitan areas received many evacuees while also experiencing a broader migration boom. These background flows complicate the interpretation of the negative slopes. The regressions provide no evidence of excess net out-migration from the main receiving cities, but they do not track individual evacuees.

Fitting $S_{\text{med}}$ and $\rho_h$ to destination shares also drives the fitted coupling toward $S_{\text{med}}\to0$. We therefore set $S_{\text{med}}=0$ in the main analysis and test values of 0.01 and 0.05 in the sensitivity analysis.

\subsubsection{Fitted dynamics}
Stage A reproduces metropolitan housing stock within 0.3 percentage points. The fitted values are $d_0=0.42$, $d_\infty=0.11$ and $\rho_s=0.76$ per year. The initial loss $d_0$ is extrapolated from the first post-storm housing estimate in July 2006. It is the least certain of these parameters, with a jackknife standard error of 0.16.

Stage B gives $S_s=1.62$, $\kappa=8.8$ and $\lambda=0.50$ per year. Fitted population is within 0.35 percentage points of the Census estimate in every year. The model values for 2006--09 are 0.740, 0.772, 0.800 and 0.812 of trend. The corresponding observations are 0.740, 0.774, 0.797 and 0.813. The fitted IRS timing share is 0.958, compared with 0.957 observed. The settling rate corresponds to a half-life of about 1.4 years.

Tracking displaced residents improves the population fit. Without tracking, the best fit has $S_s=1.59$ and $S_{\text{ret}}=0.46$. Its sum of squared population residuals is 19 times larger. The Akaike information criterion favours tracking by about 10 points, although this comparison rests on only four annual observations. The additional tests of return timing and origins provide a more informative comparison (Supplementary Section~\ref{si:validation}).

Stage C gives $\rho_h=1.9$ per year, with a wide interval of 0.4--8.6. The model accounts for 78\% of observed capacity growth across the ten receiving cities in 2006 and 59\% in 2007. Across cities, modelled and observed growth correlate at 0.45 and 0.53 in those years. Austin and San Antonio added much more housing than the model predicts from evacuee inflows, consistent with the regional housing boom. Baton Rouge added less than predicted from its occupancy.

\subsubsection{Uncertainty and sensitivity}\label{si:sensitivity}
Monte Carlo propagation of the stage B and C covariance estimates gives 95\% population bands that cover all four annual Census estimates (Fig.~\ref{fig:calib}a). These bands describe uncertainty in those fitted parameters.

The housing series is especially important for separating the initial capacity loss from the evacuation rate. We repeat stages B and C after omitting each year from the stage A housing fit. Omitting July 2006 raises $d_0$ to 0.63 and lowers $S_s$ to 0.67. Omitting any other year leaves $S_s$ between 1.43 and 1.74. The settling rate $\lambda$ remains between 0.47 and 0.62.

We also test eleven variants of network membership, baseline years, trend window, receiving-city selection, and coupling between receiving cities. The estimates of $S_s$, $\kappa$ and $\lambda$ change by less than 15\%. Validation $R^2$ values change by at most 0.04 (Table~\ref{tab:sens}). The capacity response $\rho_h$ is less stable. It ranges from 0.6 to 2.3, mainly because of changes to the trend window. A network based entirely on gravity estimates fits population equally well but predicts destinations and return origins much less well. Under that network, $\rho_h$ is not identified.

\subsection{Additional validation comparisons}\label{si:validation}
We next compare the calibrated model with additional migration and occupancy data (Fig.~\ref{fig:calib}). These outcomes were not targets for estimating evacuation, return, and settling rates. Destination shares also informed the choice of coupling between receiving cities.

\begin{figure}[tbp]
\centering
\figfile{figures/figS3_calibration.pdf}{\textwidth}
\caption{Calibration and validation. (a) New Orleans population and housing relative to the 2001--05 trend. Dots show Census population and squares show housing units. The dashed curve is capacity fitted to housing. The population curve includes a 95\% Monte Carlo band. These are stage A and B targets. (b) Net domestic in-migration relative to the 2002--05 mean. Dots show Census estimates. The solid curve tracks displaced residents and includes a 95\% band. The dashed curve uses gradient-driven return. (c) Each receiving metropolitan area's share of excess out-migration in 2005--06, compared with its share of modelled population gains. The vertical group contains metropolitan areas with suppressed flows from New Orleans, which receive the same capped weight. (d) The corresponding comparison for origins of excess in-migration in 2006--07. (e) Census net domestic migration above the 2002--05 mean in 2005--06, compared with modelled July 2006 population gains for 60 metropolitan areas within 900~km. (f) Census persons per housing unit relative to trend, compared with modelled excess occupancy in the ten main receiving cities. Filled symbols indicate July 2006 and open symbols July 2007. Panels b--f provide validation comparisons. Grey lines mark $y=x$.}
\label{fig:calib}
\end{figure}

\subsubsection{Destinations and return timing}
Across receiving metropolitan areas, the model predicts the shares of excess out-migration in 2005--06 with $R^2=0.70$. This is close to the 0.72 obtained from observed pre-disaster flow shares alone. For the origins of excess return migration in 2006--07, tracking displaced residents gives $R^2=0.55$. Gradient-driven return gives 0.32.

Modelled population gains in July 2006 also correlate with Census net domestic migration above the 2002--05 average. The correlation is $r=0.79$ across 60 metropolitan areas within 900~km. These quantities are related but differ in their coverage and definition.

The model captures the decline in Census net domestic in-migration. Predicted net in-migration to New Orleans is 43{,}000, 38{,}000 and 16{,}000 in the years ending June 2007, 2008, and 2009. The observed values are 45{,}000, 28{,}000 and 22{,}000. Gradient-driven return instead gives 23{,}000, 46{,}000 and 41{,}000, with the largest return occurring too late. For the storm year, modelled net migration is $-353{,}000$, compared with the Census estimate of $-325{,}000$ above the 2002--05 mean. The difference is within 9\%.

\subsubsection{Where the model falls short}
The model substantially underpredicts Houston's role. It assigns Houston 13\% of evacuees and 13\% of returnees, compared with observed shares of 28\% and 41\%. Houston was the main destination of organised evacuation, which increased its population by about 3\% \citep{McIntosh2008}. Pre-existing migration ties do not capture that directed movement.

The model also underpredicts receiving-city population gains by about a third. Regressing the Census measure on the model measure gives a slope of 1.55.

Return migration in the IRS data is more concentrated in the first year than the model predicts. The ratio of excess in-migration in 2006--07 to that in 2007--08 is 1.87 in the data and 0.72 in the model. Crowding in receiving cities also lasts longer in the data. Census persons per housing unit remain above trend in most main destinations in 2007, when modelled occupancy has returned to capacity. The correlation across receiving cities and years is 0.52.

\subsection{County fits and housing capacity checks}\label{si:county}
\subsubsection{Counties and data}
We next apply the model to ten counties affected by four disasters (Fig.~\ref{fig:events}). Seven were struck by Katrina. They are Orleans, St Bernard, Plaquemines, and Jefferson parishes in Louisiana, and Hancock, Harrison, and Jackson counties in Mississippi. The other cases are Cameron Parish after Rita in 2005, Aransas County in Texas after Harvey in 2017, and Butte County in California after the 2018 Camp Fire.

Each county's Census housing estimates record a loss after the disaster. By the first July after the event, that loss ranges from 3.5\% in Aransas to 78\% in St Bernard. We also examined Galveston after Ike in 2008 and Bay County, Florida, after Michael in 2018. Their Census housing estimates show no loss, so this source cannot represent their capacity shocks. They enter the larger analysis described in Supplementary Section~\ref{si:expanded}, which uses survey-based housing measures.

For these comparisons, we use the reduced county model described in Methods. It assumes that receiving areas remain near their baseline occupancy and omits local demographic change. Capacity follows each county's housing stock relative to its pre-event trend. We fit population relative to trend and the share of two-year excess IRS out-migration recorded in the event filing year. The return rate $\kappa$ is shared across counties. Evacuation and settling rates are fitted separately for each county.

\subsubsection{Housing and population losses differ across counties}
With the occupancy threshold fixed at one, the model fits Orleans, St Bernard, and Jefferson within about 1 percentage point. Errors in the other counties reach 7 points. The median root-mean-square error across counties is 2.8 percentage points (Fig.~\ref{fig:panel}c).

The discrepancies have two different directions (Fig.~\ref{fig:events}). In Hancock, Harrison, and Butte, population one year after the event exceeds the capacity implied by remaining housing at pre-event occupancy by 6--20\%. The excess then narrows. In Hancock, it falls from 20\% in 2006 to 4\% in 2009. This pattern is consistent with residents remaining locally in temporary housing or sharing homes. After Katrina, FEMA placed trailers and mobile homes on private lots and group sites across coastal Mississippi. About 134{,}000 homes there had been destroyed or damaged \citep{GAO2007,EvansCowley2011}.

Plaquemines, Cameron, and Aransas show the opposite pattern. Population falls below the capacity implied by the housing stock. In Plaquemines and Cameron, it continues to fall after housing has stabilised. Both uninhabitable homes that remain in the unit count and departures unrelated to measured housing loss could contribute to this gap.

Rates calibrated on metropolitan New Orleans produce county prediction errors of 1.2--10.3 points, with a median of 4.7. Pooling rates across the other nine county cases gives errors of 1.3--12.8 points for the omitted county, with a median of 6.4. This comparison leaves out counties, rather than whole disasters.

\subsubsection{Allowing the occupancy threshold to differ}
We then allow the occupancy threshold $\phi^*$ to vary across counties. In the reduced model, outward movement occurs above this threshold and return occurs below it. Setting $\phi^*=1$ recovers the baseline specification.

The additional parameter reduces the median fitting error to 0.5 percentage points. The largest error is 2.3 points in Cameron, where the model still misses the continuing decline. The Akaike information criterion falls from 2277 to 209. The extended fit uses 31 parameters for 47 observations. These are improvements in fit to the same population series used to estimate the thresholds. The later test on omitted disasters assesses prediction.

Estimated thresholds range from 0.80 to 1.25 (Fig.~\ref{fig:panel}a). Hancock and Harrison have values of 1.25 and 1.09. Butte and Jackson have values of 1.08 and 1.04. These values allow population to remain above the capacity implied by the housing measure. Orleans and Jefferson are close to one, at 0.97 and 1.01. Plaquemines, Aransas, and Cameron fall below one, at 0.95, 0.93 and 0.90. St Bernard reaches the lower fitting bound of 0.80, so its boundary-based interval is not meaningful.

The threshold summarises the relation between population and measured housing capacity. A value above one is compatible with local accommodation beyond that housing measure. A value below one permits population loss beyond the measured housing loss. Neither interpretation identifies a household mechanism on its own. Errors in the capacity measure can also change the fitted threshold.

Evacuation rates are identified in five counties, whose thresholds lie between 0.80 and 1.01. Across those counties, the estimated rate decreases as the housing shock increases. It is 3.5 per year in Jefferson, where 7\% of housing was lost, and 0.5 in St Bernard, where 78\% was lost. The log--log slope is $-0.77$, with $r=-0.91$ (Fig.~\ref{fig:panel}b). In the other five cases, population approaches the fitted threshold within months. The annual observations therefore do not identify the evacuation rate. The comparison of evacuation speeds rests on the five identified cases.

\subsubsection{The threshold depends on how capacity is measured}
Profile-likelihood intervals for $\phi^*$ are narrower than the 0.02 grid spacing in every county under the assumed measurement model. These intervals re-optimise evacuation and settling rates at each threshold. St Bernard's optimum remains 0.80 when the lower bound is relaxed to 0.6. Changing the capacity measure has much larger effects than these statistical intervals suggest.

Five counties have ACS coverage. Their survey-based usable housing measure excludes units vacant for other reasons, a category that includes some damaged homes. It falls farther than the Census housing-unit count. One year after Katrina, the usable housing measure is 0.41 of trend in Orleans, compared with 0.45 from the Census estimates. In Jefferson, the corresponding values are 0.83 and 0.93.

Refitting with ACS-based capacity raises the threshold in all five counties. Changes are modest in Harrison, from 1.09 to 1.11, Jackson, from 1.04 to 1.10, and Butte, from 1.08 to 1.09. They are larger in Jefferson, from 1.01 to 1.13, and Orleans, from 0.97 to 1.27. Under this capacity measure, the two New Orleans parishes also have population above the level implied by remaining usable housing.

The above-capacity pattern therefore persists across housing measures for the Mississippi counties and Butte. The classification of the large New Orleans parishes changes. We cannot make the same check for St Bernard, Plaquemines, Cameron, or Aransas because the annual ACS does not cover them (Supplementary Section~\ref{si:methods}). The survey measure is still a proxy for usable housing rather than a direct inspection of habitability.

Across the ten counties, the fitted threshold has weak correlations with pre-disaster vacancy ($r=-0.12$), non-seasonal vacancy ($r=0.11$), renter share ($r=0.12$) and shock size ($r=-0.39$). Several processes could explain the remaining differences. Nearby housing may accommodate additional households, temporary units may be added, and damaged homes may remain in the housing count. These data do not separate those explanations.

We also compare predicted second-year excess in-migration with IRS data that were not used in this fit. The Pearson correlation is $r=0.87$ across counties. It is driven by two New Orleans parishes with large returns. Removing them reduces the correlation to $-0.25$. The rank correlation across all ten counties is also weak, at 0.24.

\subsubsection{Transferring the Katrina rates to the Camp Fire}
Without refitting, the Katrina metropolitan rates predict Butte County's population at 0.90 and 0.87 of trend in July 2019 and 2020. The Census estimates are 0.94 and 0.91. Persons per housing unit rise about 10\% above trend, consistent with some households remaining within the county after losing their homes.

Using pre-disaster migration shares to predict Butte's destinations gives $R^2=0.21$. The data identify 68 counties with positive excess arrivals. Sacramento receives 4\% of excess out-migration, compared with an 11\% pre-disaster share. Fitting Butte's own rates and occupancy threshold, estimated at 1.08, reduces population error from 4.3 to 1.6 percentage points.





\subsection{Expanded sample and prediction diagnostics}\label{si:expanded}
\subsubsection{Extending the sample with survey-based housing data}
The Census housing-unit estimates do not capture a loss after every disaster. The ACS provides another measure. For counties with at least 65{,}000 residents, its annual estimates distinguish occupied and available units from units vacant for other reasons. The latter category includes damaged homes awaiting repair. We define usable housing as total units minus other-vacant and seasonal units, then express this measure relative to its pre-event ACS trend.

We listed 42 county--disaster pairs before fitting and later added an omitted case, Orleans Parish after Ida, bringing the total to 43. The list covers counties with ACS data named in major housing-damaging disasters from 2008 to 2024. These include hurricanes Ike, Sandy, Matthew, Harvey, Irma, Florence, Michael, Laura, Ida, Ian, and Helene. They also include the Tuscaloosa, Joplin, and Moore tornadoes, the 2016 Louisiana floods, and the Tubbs, Carr, Camp, and Lahaina fires. All listed cases were analysed regardless of measured housing loss. Thirty-seven have enough ACS years for the analysis.

Most measured first-year losses are small. Butte loses 15\% of usable housing after the Camp Fire and already belongs to the Census panel. Eleven additional counties lose at least 4\%. The largest losses are 13\% in Calcasieu after Laura, 12\% in Monroe after Irma and 10\% in Bay after Michael. In eleven cases, the ACS records no loss. Modest county-wide damage and survey sampling error of about 1--2\% of units both affect these measurements.

Among the eleven additional counties with at least 4\% loss, eight have fitted thresholds above one. The median is 1.03. Identified values range from 0.92 to 1.07, while Galveston's threshold is not identified. In each of these counties, population one year after the event exceeds the capacity implied by remaining usable housing at pre-event occupancy by 2--9\%. Average household size provides an additional check on household sharing. It was not used in fitting and does not track the threshold across counties. The correlation is $r=0.08$, or $-0.20$ without Galveston.

\subsubsection{Leaving out one disaster at a time}
We test whether parameters estimated from other disasters predict population change in a new case. The test combines the ten Census-panel counties with the eleven additional ACS counties. This gives 21 counties affected by 14 disasters, with the Camp Fire counted once.

We fit two versions of the model. One fixes $\phi^*=1$. The other uses partial pooling, which pulls each county's threshold towards a common estimated mean. The assumed standard deviation around that mean is 0.1. For each disaster in turn, we fit both models to all other disasters and predict every county affected by the omitted event.

The omitted counties use their measured post-event capacity paths. Their evacuation rates come from the relation between housing loss and evacuation rate in the training data. The pooled model uses the training mean for $\phi^*$. No post-event population observations from the omitted disaster enter these predictions. This is a test of population prediction conditional on observed housing loss and recovery. It does not forecast the capacity path before the disaster occurs.

The pooled mean threshold is 1.03 in the full fit and ranges from 1.02 to 1.03 across training sets. We measure each county's error as the root-mean-square difference between predicted and observed population relative to trend. The median is 3.0 percentage points, with county errors ranging from 0.8 to 10.4 (Fig.~\ref{fig:holdout}). The main text compares the alternatives in Table~\ref{tab:heldout_summary}.



The pooled model improves on population proportional to capacity in 16 of 21 counties and on the $\phi^*=1$ model in 14. The $\phi^*=1$ model also improves on proportional capacity tracking in 16 counties. Much of the gain over that baseline therefore comes from modelling adjustment over time. For the ten counties with larger housing losses, the median paired reduction relative to capacity tracking is 2.2 points. Its 95\% interval ranges from $-0.2$ to 4.8, so this smaller comparison remains uncertain.

In St Bernard and Orleans, the pooled model follows a slower population recovery than capacity tracking predicts. Its largest errors occur where a common threshold poorly represents the county. The error is about 10 points in Hancock and 6--9 points in Cameron, Plaquemines, and Aransas. For small events, population changes little relative to trend. In those cases, assuming no departure from trend performs about as well as the dynamic model.

\subsubsection{Sensitivity of the prediction results}
We quantify uncertainty by resampling whole disasters in a cluster bootstrap with 10{,}000 draws (2{,}000 for the sensitivity variants). The bootstrap keeps counties from the same disaster together. The county-level Wilcoxon tests do not account for this dependence, so we place more weight on the bootstrap intervals. Relative to proportional capacity tracking, the median paired error reduction is 1.9 percentage points, with a 95\% interval of 0.4--2.5. A county-level Wilcoxon signed-rank test gives $p=0.006$. The paired reduction remains 1.8--2.0 points when the threshold standard deviation is set to 0.05 or 0.2 instead of 0.1.

Changing the ACS inclusion threshold also changes the sample. A 3\% housing-loss threshold gives 24 counties and a paired reduction of 1.8 points. A 5\% threshold gives 17 counties and a reduction of 1.4 points, with $p=0.19$. The gain over the fixed-threshold dynamic model is smaller. Its median paired reduction is 0.5 points, with an interval of 0.2--1.3 and $p=0.17$. Across these variants, the reduction is 0.5--0.7 points.

We repeat the analysis with ACS-based capacity for the five Census-panel counties that have survey coverage (Supplementary Section~\ref{si:county}). The pooled threshold rises to 1.06, while median error remains 3.0 points. Capacity tracking has a median error of 6.1 points. The median paired reduction is 2.1 points, with an interval of 0.3--3.5 and $p=0.03$. However, errors increase for all four Katrina counties in this subset. In Orleans, error rises from 4.6 to 13.3 points. A small change in the pooled threshold therefore does not improve every county.

The dynamic model also improves on three statistical regressions and two partial-adjustment forecasts fitted to the same training disasters. Regressing the population gap on the capacity gap gives a median error of 5.6 points. Adding an interaction between capacity gap and time since the event gives 5.2 points. Adding the first-year housing loss gives 4.0 points. A partial-adjustment forecast with a fitted rate gives 4.6 points, with a median fitted rate of 2.5 per year across training sets. Letting the long-run population gap be a fitted multiple of the capacity gap gives 6.0 points. The dynamic model's median is 3.0 points.

The dynamic model has lower error in 15, 16, 18, 15, and 19 of the 21 counties for these five comparisons. Wilcoxon $p$ values are 0.010, 0.018, 0.001, 0.006, and 0.001. Cluster-bootstrap intervals for the difference in median error are 1.4--4.0, 1.4--3.5, 0.6--2.9, 0.8--2.8, and 1.3--3.6 points, respectively.

With the threshold fixed at one, the dynamic model's median error is 3.6 points. It has lower error than the five forecasts in 14, 15, 17, 15, and 18 counties. Cluster-bootstrap intervals for the difference in median error are 0.2--3.5, $-0.02$--3.1, $-0.4$--1.7, $-0.4$--2.3, and 0.2--3.4 points. A clear advantage over the stronger forecasts therefore requires the pooled threshold.

Prediction intervals include variation in occupancy thresholds and evacuation rates across counties. The 90\% intervals average 19 percentage points in width. They cover 95\% of the 80 omitted county-year observations. The 50\% intervals cover 71\%. Both exceed their nominal coverage in this evaluation. These intervals hold the measured housing path and shared parameters fixed, so they do not represent uncertainty in the full prediction process.

We also predict with less housing information for the omitted county (Methods). With the first post-event loss and the median training rebuilding curve, the pooled model's median error is 2.5 points, compared with 3.0 using the measured path. Capacity tracking with the same curve gives 4.5 points. With no rebuilding assumed, the errors are 2.8 and 4.4 points. For the eleven counties with losses below 10\%, the typical curve reduces the model's median error from 1.9 to 1.4 points. It probably smooths sampling noise in survey estimates of small losses. For the ten counties with larger losses, median errors are 3.6 with the measured path, 4.2 with the typical curve, and 5.3 with no rebuilding. Capacity tracking with the same inputs gives 7.7, 8.3, and 7.0.



\subsection{Hazard geography}\label{si:hazard}
We test whether nearby housing outside hazard zones helps explain differences in the occupancy threshold. The hypothesis is that counties with more housing outside the affected zone can accommodate more households locally. We examine the 17 counties affected by ten storms or floods in the prediction sample. We could not obtain comparable hazard maps for the three fires and the tornado.

For every 2010 Census block, we record its elevation and whether it lies in a FEMA Special Flood Hazard Area, a NOAA Category~3 storm-surge zone, or a NOAA levee area. Levee areas are protected low ground that the surge model does not resolve. We calculate the share of pre-disaster housing outside all three zones. For the two inland floods, the 2016 Louisiana floods and Matthew's flooding in Robeson County, we use the flood zone alone. For coastal storms, we also calculate the housing share at least 3~m above sea level.

We fixed these definitions before calculating shares for the full sample. A preliminary check in St Bernard showed that the current flood maps omit much leveed land, which motivated including levee areas. We also report the flood-zone measure separately. The housing distribution refers to pre-event periods, but hazard zones come from current map products. This is therefore a retrospective comparison, rather than a test using the maps available before each event.

The fitted thresholds broadly follow the hypothesised ordering (Fig.~\ref{fig:hazard}a,b). Almost all housing in the four Louisiana parishes struck by Katrina and in Cameron lies inside the mapped hazard area and below 3~m. Their thresholds range from 0.80 to 1.01, with St Bernard at the fitting bound. In Mississippi, Hancock, Harrison, and Jackson have 28--44\% of housing outside the mapped hazard area and 83--97\% above 3~m. Their thresholds range from 1.04 to 1.25.

Across the 17 counties, the housing share outside the hazard area has a Spearman correlation of $\rho=0.55$ with the threshold, with $p=0.02$. The Pearson correlation is $r=0.40$, with $p=0.11$. For housing above 3~m in the 15 coastal counties, the corresponding values are $\rho=0.68$ with $p=0.005$ and $r=0.45$ with $p=0.09$. The flood-zone measure alone has a much weaker rank correlation, $\rho=0.15$. Current FEMA flood maps exclude most of the leveed New Orleans area that flooded in 2005.

The association is strongest within Katrina, with $\rho=0.87$ across its seven counties. Across the other disasters, $\rho=0.60$, $p=0.07$ and $n=10$. Monroe County after Irma is an exception to the expected pattern. It is almost entirely low-lying but has a fitted threshold of 1.06.

Several features make these associations uncertain. Four counties have a hazard share of zero and five have shares below 1\%. Pearson correlations are not significant. We examined several measures without adjusting for multiple comparisons. Replacing the capacity measure with ACS-based estimates in Orleans, Jefferson, Harrison, and Jackson largely removes the association. The rank correlation becomes 0.06 for the hazard-area measure and 0.36 for elevation.

The threshold estimates also come from two model fits with slightly different bounds. They include Galveston, whose threshold is not identified. For Jefferson County, Texas, the flood layer contains no Special Flood Hazard Areas. Its hazard measure therefore relies entirely on surge and levee areas.

FEMA assistance registrations provide a separate, post-event measure of where damage occurred. Across 20 counties, we calculate the housing share in ZIP codes with fewer than ten registrations per 100 units. Its rank correlation with the threshold is $-0.44$, opposite to the expected direction. Thresholds of 20 and 30 registrations give correlations of $-0.41$ and $-0.36$. Registrations may also reflect assistance that helps households remain locally, which complicates their interpretation as a damage measure.

We then ask whether geography improves prediction for omitted disasters. For the 17 storm and flood counties, the mean occupancy threshold is fitted as a linear function of the housing share outside the hazard area. Median population error is 4.1 points, compared with 3.7 for a common pooled threshold. Geography improves prediction in eight of the 17 counties (Fig.~\ref{fig:hazard}c). With the elevation measure, the median errors are 3.9 and 4.4 points, and geography improves prediction in nine of 15 counties.

The fitted hazard-area effect is small. Depending on the training set, the threshold increases by at most 0.06 across the observed range of housing shares. In the joint fit, the Louisiana parish thresholds that drive much of the contrast are weakly identified and pulled towards the common mean. When Katrina is omitted, the remaining disasters show only a weak relation. Elevation gives a modest improvement, but the result does not hold consistently across geographic measures or capacity definitions.

\begin{figure}[tbp]
\centering
\figfile{figures/figS4_hazard.pdf}{\textwidth}
\caption{Occupancy thresholds and mapped hazard geography. (a) Fitted threshold against the share of pre-disaster housing outside FEMA flood zones, NOAA Category~3 surge zones, and NOAA levee areas. Inland flood cases use the flood zone alone. (b) Threshold against the share of housing at least 3~m above sea level for coastal storms. Filled symbols show Katrina counties and open symbols show other disasters. Bars are 95\% intervals. Galveston has no interval because its threshold is not identified. (c) Population prediction error when the threshold depends on the hazard measure in panel a, compared with error under a common pooled threshold. Both axes are logarithmic. Filled symbols indicate Katrina counties and the dashed line marks equal error. Hazard zones use current map products.}
\label{fig:hazard}
\end{figure}

\subsection{Temporary accommodation and reported time away}\label{si:accommodation}
Two public sources provide further evidence on displacement and accommodation, although neither covers the full sample. FEMA registrant records cover six panel counties affected by Harvey, Irma, Florence, Michael, and Laura. They do not record the location of housing rented with assistance, so we cannot determine whether those households remained in the county.

The records do identify stays in FEMA-paid hotels under Transitional Sheltering Assistance. The hotel-placement share is lower where the fitted occupancy threshold is higher. Spearman's $\rho$ is $-0.94$ across six counties and $-0.90$ after excluding Laura, for which no hotel stays are recorded. This small comparison is compatible with differences in accommodation options. It does not establish where displaced households stayed or identify the mechanism behind the fitted threshold.

From December 2022, the Census Household Pulse Survey asked adults whether a disaster had displaced them in the preceding year. In Hawaii after the Lahaina fire, 1.6\% of adults reported fire displacement lasting more than a month. This represents about 15{,}000 adults, with a standard error of 0.25 percentage points. The estimate is about twice the approximately 7{,}000 adults associated with the usable housing loss measured by the ACS on Maui. Among all adults displaced by fire, including those displaced for less than a week, 36\% had not returned home when surveyed.

In Florida after Ian, 1.1\% of adults reported displacement lasting more than a month, representing about 155{,}000 people. This is six times the estimated number of adults associated with housing lost in the three affected counties with ACS coverage. The geographic coverage differs. Household Pulse covers the whole state, while Ian and Nicole, which followed six weeks later, damaged homes in many counties.

These comparisons indicate that housing-unit losses and reported displacement measure different outcomes over different areas. Short-term loss of access to a home may occur without a corresponding reduction in the annual housing measure. This helps explain why the fitted occupancy threshold can reflect limitations of the capacity data as well as local accommodation.

\subsection{Cumulative excess occupancy}\label{si:sim_measure}
For simulations of the full network, we measure the number of residents above capacity and how long that excess lasts,
\begin{equation}
Q(T)=\frac{1}{N_s(0)}\int_0^T\sum_i(N_i-K_i)_+\,dt.
\label{eq:severity}
\end{equation}
The sum includes the affected area and every host. The measure has units of excess-occupancy person-years per original resident of the affected area. Population is unchanged at the instant of the shock, so $N_s(0)$ is also its pre-disaster population.

This is a model measure of cumulative pressure on housing. We refer to it as severity in the simulation analysis. Relating it to observed crowding requires a consistent definition of capacity and occupancy data.

The current specification replaces the prototype developed at the Santa Fe Institute Complex Systems Summer School. Supplementary Section~\ref{si:proto} explains the changes and the problems identified in the earlier implementation.

\subsection{Limiting cases for cumulative excess occupancy}\label{sec:theory}
Consider a limiting case in which growth is slow ($r\ll S_s$) and the affected city does not rebuild. Receiving cities add enough capacity to keep their occupancy near one as people arrive. The affected city's occupancy then approximately follows $d\phi_s/dt\approx-S_s\phi_s(\phi_s-1)$, starting from $\phi_s(0)=1/(1-d_0)$. Under this approximation, its cumulative excess occupancy is
\begin{equation}
\int_0^\infty (N_s-K_s)\,dt = \frac{K_s}{S_s}\ln\frac{1}{1-d_0},
\label{eq:regime1}
\end{equation}
Severity therefore scales as $S_s^{-1}$. Evaluating population loss through the growth term along the same trajectory gives $r\,d_0/S_s$ per resident. In this regime, faster evacuation reduces crowding because receiving cities have room for arrivals.

If evacuation is fast relative to capacity expansion in receiving cities, crowding shifts to those destinations. Severity can then be limited by their response rate, $\rho_h$. Supplementary Sections~\ref{si:calibration} and~\ref{sec:sims} examine these regimes using the Katrina calibration.

\subsection{Simulations of cumulative excess occupancy}\label{sec:sims}
We apply the calibrated Katrina shock and rates to each of the 196 metropolitan areas in turn. Destinations follow the observed migration network (Fig.~\ref{fig:sims}). We measure cumulative excess occupancy over the first six years after the shock, using $Q(6)$. With a common evacuation rate, severity changes little with the location of the shock. It ranges from 0.239 to 0.250 excess-occupancy person-years per resident, compared with 0.242 for New Orleans. Receiving cities account for half of this total.

In these simulations, receiving cities do not pass evacuees on. They add capacity at a rate proportional to their own crowding. The identity and network position of the destinations therefore have little effect on total severity under these assumptions.

We also test a scenario in which the evacuation rate is proportional to the affected city's pre-disaster per-capita out-migration rate. This is an assumption about evacuation behaviour. In this scenario, severity decreases with the evacuation rate, with a log--log slope of $-0.26$. The decline is weaker than the slope of $-1$ in Eq.~\eqref{eq:regime1}, because receiving cities' contribution does not depend on the evacuation rate.

Severity rises smoothly with initial capacity loss, with no threshold in the simulated range. It is much more sensitive to the receiving cities' capacity response. Severity falls from 1.14 to 0.15 person-years per resident as $\rho_h$ increases from 0.05 to 10 per year. Across the calibrated 95\% interval for $\rho_h$, severity varies by a factor of three to four.

The prototype linked network position to severity through onward migration from crowded receiving cities. That mechanism is absent from the calibrated model, which sets onward coupling to zero. The available aggregate migration data support this modelling choice but cannot rule out onward movement by individual evacuees. Among the parameters examined, the receiving cities' capacity response has the largest effect on severity and is also the least precisely estimated.

\begin{figure}[tbp]
\centering
\figfile{figures/figS5_simulations.pdf}{\textwidth}
\caption{Severity in the calibrated model over the first six years after the shock, measured in excess-occupancy person-years per resident of the affected city. (a) Severity when each of the 196 metropolitan areas is affected, plotted against its evacuation rate. Circles use rates proportional to pre-disaster out-migration intensity. The filled circle marks New Orleans. The grey band shows the range when all cities use the calibrated New Orleans rate. (b) Severity after a New Orleans shock as a function of capacity lost at landfall. The solid line shows total severity, and the dashed line shows the contribution from receiving cities. The dotted line marks the calibrated capacity loss. (c) The same outcomes as a function of the receiving cities' capacity response. The grey band marks its calibrated 95\% interval.}
\label{fig:sims}
\end{figure}

\subsection{Relation to the prototype}\label{si:proto}
An earlier version was developed at the Santa Fe Institute Complex Systems Summer School. It used a single coupling constant, with network weights normalised by the strongest connection. Receiving cities expanded toward $1.3K_i$ at a fixed build rate. Expansion began at the time of the disaster, whether or not those cities were crowded. Distances were measured in degrees of latitude and longitude, and displaced residents were not tracked separately.

A code audit showed that the prototype produced population growth without a disaster and overstated long-run population loss in the affected city. The results in this paper replace those from the prototype. The repository contains the code, tests, and audit, as described in Data and code availability.

\subsection{Supplementary methods}\label{si:methods}
\paragraph{Gravity estimation.} We fit $E[y]=\exp(c+\alpha\log P_o+\beta\log P_d-\gamma\log d)$ to all ordered pairs of the 196 metropolitan areas. Here $y$ is the mean flow in 2003--04 and 2004--05. Estimation uses Poisson pseudo-maximum likelihood with HC1 standard errors. We also use a cluster bootstrap over metropolitan areas with 200 replicates.

For New Orleans, we fit $E[y_j]=\exp(c+\beta\log P_j-\gamma\log d_j)$ across destination or origin counties and metropolitan areas. Excess flows equal observed flows minus their 2003--05 mean, with negative values set to zero.

\paragraph{Calibration.} We integrate the model with LSODA, using a relative tolerance of $10^{-6}$. Stage A fits the housing path by nonlinear least squares. Its standard errors come from a leave-one-year-out jackknife. Stages B and C use weighted nonlinear least squares on log-transformed parameters. The residual variance for each target group is estimated from the fit and updated once. Parameter covariance uses the Gauss--Newton approximation.

We alternate stages B and C four times. After the second pass, $(S_s,\kappa,\lambda)$ are unchanged to three decimal places. Solver tolerance causes $\rho_h$ to vary between 1.90 and 1.93. On the log scale, this variation is about 2\% of its standard error.

Model flows are accumulated to April of each year for comparison with IRS filing years and to 30 June for comparison with Census components of change. Receiving-city housing targets are units above the 2001--05 linear trend, multiplied by the city's 2005 persons per unit.

\paragraph{Uncertainty and sensitivity.} Monte Carlo propagation uses 200 draws of $(S_s,\kappa,\lambda)$ from their joint log-scale Gauss--Newton covariance. We also draw $\rho_h$ from the lognormal distribution represented by its interval. For the jackknife analysis, we refit stages B and C for each leave-one-year-out stage-A fit.

The sensitivity analysis reruns the full pipeline under alternative network and data choices. We vary the metropolitan population threshold to 150{,}000 or 500{,}000 residents and the radius around New Orleans to 300 or 900~km. We also use a single baseline year, trend windows of 2002--05 and 2003--05, and receiving-city target sets of 5 or 20 metropolitan areas. Finally, we test $S_{\text{med}}=0.01$ and 0.05 and a network constructed entirely from gravity weights.

\paragraph{Ten-county panel.} For each county, population and housing units are expressed relative to linear trends fitted to five pre-event 1 July estimates. These years are 2001--05 for the 2005 disasters, 2013--17 for Harvey, and 2014--18 for the Camp Fire. Post-event observations extend to July 2009 for the 2005 disasters and to July 2020 for the later events.

When at least four post-event housing estimates are available, we fit Eq.~\eqref{eq:Ks} to the housing ratios. Otherwise, capacity drops at the event to the first observed ratio and changes linearly between subsequent estimates.

IRS US-total records provide out-migration totals that include suppressed flows. The timing target is the event-year share of two-year excess out-migration, $e_1/(e_1+\max(e_2,0))$. Here $e_1$ and $e_2$ are excess flows in the event filing year and the following year, measured against the mean of the two preceding years.

We minimise weighted squared errors using measurement standard deviations of 0.005 for population ratios and 0.03 for the timing share. The return parameter $\kappa$ is shared across the panel. Each affected county has its own $(S_s,\lambda)$ and, in the extended model, its own occupancy threshold $\phi^*\in[0.8,1.4]$. Parameter intervals use the Gauss--Newton covariance. Pre-disaster vacancy and tenure come from the 2000 and 2010 Census demographic profiles.

\paragraph{Larger panel.} Events, counties, and dates were listed before fitting in \texttt{13\_acs\_panel.py}. We define usable housing as ACS housing units minus seasonal units and units vacant for other reasons. This is a proxy for housing capacity. We express it relative to a linear trend fitted to at least three of the six ACS years before the event. The event year is excluded, and each ACS year $y$ is placed at mid-year.

Population is measured relative to a trend fitted to the five preceding 1 July estimates. We use intercensal estimates for 2000--09, Vintage 2020 for 2010--19, and Vintage 2025 for 2020--25. The last series is rescaled to match Vintage 2020 at 2020. Observation windows extend for 4.6 years or end when the county experiences another disaster, whichever comes first.

Fitting follows the ten-county analysis, with two changes. Both $\kappa$ and $\lambda$ are shared across the panel, and the occupancy threshold is restricted to $\phi^*\in[0.7,1.5]$. The prediction tests use the same bounds.

\paragraph{Damage geography.} For each event, we divide FEMA valid registrations from owners and renters in each ZIP code by housing units in the corresponding ZIP code tabulation area. We use the 2000 gazetteer for 2005 events and the 2010 gazetteer for later events. Areas with fewer than 100 housing units are excluded.

\paragraph{Hazard geography.} For each county, we locate the internal point of every 2010 Census block. We test whether each point falls within a FEMA Special Flood Hazard Area polygon, NOAA's Category~3 surge zone, or a NOAA levee area. The surge map is read at zoom level 14, corresponding to about 8--9~m at these latitudes. A coloured pixel indicates inundation, and a grey hatched pixel indicates a levee area. Terrain tiles provide elevation at the same resolution.

For events from 2010 onward, shares are weighted by 2010 block housing. For 2005--2008 events, Census 2000 block housing is first allocated to 2010 tracts according to intersecting land area. Within each tract, it is then distributed in proportion to 2010 housing.

The calculations run in a web browser because the analysis server cannot access the map services. The repository contains the script \path{browser/hazard_geography.js} and its output \path{data/hazard_geography.json}. In the prediction test that includes geography, the prior mean of $\phi^*_j$ is $a+bx_j$. Both $a$ and $b$ are estimated, and $\tau=0.1$.

\paragraph{Held-out prediction.} Model H allows the occupancy threshold to vary by county and adds the penalty $\sum_j(\phi^*_j-\mu)^2/\tau^2$. We set $\tau=0.1$ and estimate $\mu$. The parameters $\kappa$ and $\lambda$ are shared, while $S_s$ and $\phi^*$ vary by county. Model A fixes $\phi^*=1$.

We omit one disaster at a time and refit both models to the counties affected by the remaining disasters. Across these training counties, we regress $\log S_s$ on log housing loss. Each omitted county is then simulated using its measured capacity path and the evacuation rate given by this regression. Its occupancy threshold is set to $\mu$ for model H and to 1 for model A. These predictions therefore use observed post-event housing capacity, while population data from the omitted disaster are not used for fitting.

We repeat the analysis with $\tau=0.05$ and 0.2 and with survey inclusion thresholds of 3\% and 5\%. Uncertainty intervals use a cluster bootstrap over omitted disasters. We also apply paired Wilcoxon signed-rank tests to county errors. These analyses are implemented in \path{20_holdout_uncertainty.py}.

\paragraph{Capacity check and profile likelihood.} The ACS provides no pre-event year for Orleans, Jefferson, Harrison, or Jackson. For these counties, we construct the reference housing series from the 2000--05 linear trend of Census housing estimates. We multiply this trend by the Census 2000 share of usable units, defined as one minus the seasonal and other-vacant shares in SF1 table H5. ACS usable housing is divided by this reference series. Butte uses the ACS capacity measure from the larger panel.

We fit each county separately, fixing $\kappa$ at its ten-county estimate. The parameters $S_s$, $\lambda$, and $\phi^*\in[0.6,1.6]$ remain free. Fits start from the panel estimates and from five additional starting points. To construct the profiles, we fix $\phi^*$ on a grid from 0.70 to 1.50 in steps of 0.02. At each value, we re-optimise $S_s$ and $\lambda$, continuing outward from the optimum. The implementation is in \texttt{23\_capacity\_check.py}.

\paragraph{Benchmarks and prediction intervals.} The statistical forecasts use the same omitted-disaster groups as the dynamical model. Each regression uses measured capacity at the observation time and has no intercept. The first fits $N-1=b(K-1)$. The second adds $(K-1)t$, and the third also includes the first-year capacity gap $K_0-1$. The partial-adjustment forecasts start with no population gap at the event and update at each observation time; the second variant replaces the target $g_k$ by $bg_k$. Rates and $b$ are fitted by nonlinear least squares to the training counties. These baselines are implemented in \texttt{24\_baseline.py}.

Prediction intervals use 400 draws per county. We draw $\phi^*\sim N(\mu,0.1^2)$ and clip values outside $[0.7,1.5]$ to the nearest bound. We draw $\log S_s$ from the training regression on log housing loss, using its residual standard deviation. The measured housing path, shared rates, and regression coefficients remain fixed across draws. The implementation is in \texttt{26\_pred\_intervals.py}.

\paragraph{Destinations beyond Katrina.} For each destination county, excess out-migration equals the IRS flow in exemptions during the event filing year minus the mean of the two preceding years. Negative values and flows below the disclosure threshold are set to zero.

We compare observed destination shares with three predictions. These use the affected county's mean baseline flow shares, a gravity kernel with the national exponents, or a kernel with the New Orleans county-level exponents. The national size and distance exponents are 0.90 and 1.12. The New Orleans exponents are 1.42 and 2.16. Here $R^2$ denotes the squared correlation across destinations. The analysis is implemented in \texttt{25\_destinations.py}.

\paragraph{FEMA registrants and Household Pulse.} We classify primary-residence registrants as displaced if their home was destroyed or required habitability repairs. Hotel stays count registrants who checked in under Transitional Sheltering Assistance. The rental-resource location field is blank in almost all records.

Household Pulse estimates use person weights and 80 replicate weights. Variances use successive differences with a factor of 4/80. Week 55 could not be downloaded and is omitted. To estimate the number of adults in lost housing, we multiply ACS usable units lost by persons per household and by 0.78, the adult share of household members. These calculations are in \texttt{21\_fema\_mechanism.py} and \texttt{22\_pulse.py}.

\clearpage
\subsection{Supplementary tables}
\begin{table}[h]
\centering\small
\caption{Gravity estimates for flows out of and into New Orleans, fitted by Poisson pseudo-maximum likelihood with robust standard errors. The county analysis includes 683 counties. These comprise every county with a reported flow and all other counties with at least 100{,}000 residents. $N$ is the number of exemptions modelled.}
\label{tab:gravity}
\begin{tabular}{lcccc}
\toprule
Flow & $\hat\beta$ (population) & $\hat\gamma$ (distance) & $N$ & pseudo-$R^2$\\
\midrule
Out, baseline 2003--05 & $1.42\ (0.06)$ & $2.16\ (0.09)$ & 19{,}761 & 0.79\\
Out, 2005--06 & $1.55\ (0.13)$ & $2.08\ (0.12)$ & 183{,}052 & 0.83\\
Excess out, 2005--06 & $1.57\ (0.15)$ & $2.07\ (0.14)$ & 163{,}292 & 0.81\\
Out, 2006--07 & $1.56\ (0.07)$ & $2.34\ (0.11)$ & 30{,}561 & 0.82\\
In, baseline 2003--05 & $1.40\ (0.08)$ & $2.02\ (0.09)$ & 15{,}441 & 0.77\\
In, 2006--07 (return) & $1.78\ (0.12)$ & $2.30\ (0.12)$ & 47{,}870 & 0.86\\
\bottomrule
\end{tabular}
\end{table}

\begin{table}[h]
\centering\small
\caption{Sensitivity of calibrated parameters and validation statistics to 11 alternative specifications. Results for a network constructed entirely from gravity weights are reported separately.}
\label{tab:sens}
\begin{tabular}{lccc}
\toprule
Quantity & Main & Range over 11 variants & Gravity network\\
\midrule
$S_s$ (yr$^{-1}$) & 1.62 & 1.61--1.64 & 1.57\\
$\kappa$ (yr$^{-1}$) & 8.8 & 8.7--9.8 & 9.0\\
$\lambda$ (yr$^{-1}$) & 0.50 & 0.50--0.51 & 0.51\\
$\rho_h$ (yr$^{-1}$) & 1.9 & 0.6--2.3 & not identified\\
Destinations, $R^2$ & 0.70 & 0.67--0.71 & 0.35\\
Return origins, $R^2$ & 0.55 & 0.54--0.57 & 0.25\\
Host domestic migration, correlation & 0.79 & 0.78--0.81 & 0.84\\
Host domestic migration, slope (data on model) & 1.55 & 1.45--1.63 & 3.89\\
\bottomrule
\end{tabular}
\end{table}

\end{document}
